\section{Group Actions}
\label{sec:group-actions}

We return to the aspect where we started, that groups describe symmetries. Groups become groups of symmetry when they \textbf{act}.

\subsection{Group Actions and Homomorphisms}
\label{sec:group-actions-homomorphisms}

\begin{definition}[Group Action]
    \label{def:group-action}%
    An \emph{action} of a group \(G\) on a set \(X\) is a map
    \[
        \function{\alpha}{G \times X}{X}{\pqty{g, x}}{\alpha\pqty{g, x} = gx}
    \]
    satisfying
    \begin{enumerate}
        \item \(\alpha\pqty{e, x} = ex = x\) for all \(x \in X\): the identity acts as an identity;
        \item \(\alpha\pqty{g \cdot h, x} = \alpha\pqty{g, \alpha\pqty{h, x}}\), or \(\pqty{g \cdot h}x = g \pqty{hx}\) for all \(g, h \in G\) and \(x \in X\): composition is compatible.
    \end{enumerate}

    We write \(G \actson X\) to indicate that \(G\) acts on \(X\).
\end{definition}

\begin{examples*}[Examples of Group Actions]
    \begin{enumerate}
        \item For any group \(G\) and set \(X\), the action
        \[
            gx = x
        \]
        is the \emph{trivial action}.
        \item \(\Sym\pqty{X}\) acts on \(X\) for any set \(X\) by
        \[
            fx = f\pqty{x}.
        \]
        \item If \(G \actson X\) and \(H \subgroup G\), then \(H \actson X\) by restriction.
        \item In particular, \(\Isom\pqty{\CC} \subgroup \Sym\pqty{\CC}\) acts on \(\CC\).
        \item Similarly, \(D_{2n}\) acts on \(X_n\), the regular \(n\)-gon. It acts on the set of vertices
        \[
            \set{z \in \CC}{z^n = 1}.
        \]
        \item Every group \(G\) acts on \emph{itself} by left-multiplication: let \(X = G\), and \(G \actson G\) by
        \[
            g\gamma = g \cdot \gamma
        \]
        for \(g \in G\), \(\gamma \in X = G\). This is the \emph{left regular action}.
    \end{enumerate}
\end{examples*}

We may think of group actions in another useful way.

\begin{theorem}[Actions Define Homomorphisms]
    \label{thm:action-homomorphism}%
    An action of a group \(G\) on a set \(X\) is the same as a homomorphism
    \[
        \functiondomain{\phi}{G}{\Sym\pqty{X}}.
    \]
\end{theorem}

\begin{proof}
    Suppose \(G \actson X\). Consider the map
    \[
        \function{t_g}{X}{X}{x}{gx}
    \]
    for any \(g \in G\). Now,
    \begin{align*}
        t_{g\inverse} \pqty{t_g x} &= t_{g\inverse} \pqty{gx} \\
        &= g\inverse \pqty{gx} \\
        &= \pqty{g\inverse g} x \\
        &= ex \\
        &= x
    \end{align*}
    so \(t_{g\inverse} \circ t_g = \identity_{X}\).

    Therefore, \(t_{g\inverse}\) is \(\pqty{t_g}\inverse\) and hence \(t_g\) is a bijection.

    This tells us that \(t_g\) is a permutation on \(X\),
    \[
        t_g \in \Sym\pqty{X}.
    \]

    Therefore, we define
    \[
        \function{\phi}{G}{\Sym\pqty{X}}{g}{t_{g},}
    \]
    and we prove that \(\phi\) is a homomorphism. Take \(g, h \in G\), we have
    \[
        \phi\pqty{g} \phi\pqty{h} = t_g \circ t_h,
    \]
    so for all \(x \in X\),
    \begin{align*}
        \phi\pqty{g} \phi\pqty{h} \pqty{x} &= t_g \circ t_h \pqty{x} \\
        &= t_g \pqty{t_h \pqty{x}} \\
        &= t_g \pqty{hx} \\
        &= g \pqty{hx} \\
        &= \pqty{gh} x \\
        &= t_{gh} \pqty{x} \\
        &= \phi\pqty{gh} \pqty{x},
    \end{align*}
    and hence, \(\phi\pqty{g} \phi\pqty{h} = \phi\pqty{gh}\).

    Therefore, \(\phi\) is a homomorphism as required.

    Conversely, given a homomorphism
    \[
        \functiondomain{\phi}{G}{\Sym\pqty{X}}
    \]
    we may define an action \(G \actson X\) by \(gx = \phi\pqty{g} x\). We check this is indeed an action:
    \begin{enumerate}
        \item Identity: \(ex = \phi\pqty{e} \pqty{x} = \identity_{X} \pqty{x} = x\).
        \item Compatibility:
        \begin{align*}
            \pqty{gh} x &= \phi\pqty{gh} \pqty{x} \\
            &= \phi\pqty{g} \phi\pqty{h} \pqty{x} \\
            &= \phi\pqty{g} \pqty{\phi\pqty{h} \pqty{x}} \\
            &= \phi\pqty{g} \pqty{hx} \\
            &= g \pqty{hx}.
        \end{align*}
    \end{enumerate}
\end{proof}

The left regular action allows us to think of groups as subgroups of symmetric groups.

\begin{theorem}[Cayley's Theorem]
    \label{thm:cayley}%
    Every group \(G\) is isomorphic to a subgroup to some \(\Sym\pqty{X}\). Furthermore, if \(\abs{G} < \infty\), we may choose \(X\) to be \(\abs{X} < \infty\), or even stronger, \(\abs{X} \leq \abs{G}\).
\end{theorem}

\begin{proof}
    We choose in particular \(X = G\), and the left regular action gives \(G \actson X = G\). By \autoref{thm:action-homomorphism}, this is equivalent to a homomorphism
    \[
        \functiondomain{\phi}{G}{\Sym\pqty{X}}
    \]

    Let \(H = \img \pqty{\phi} \subgroup \Sym\pqty{X}\) be a subgroup. Then we may think of \(\phi\) as a \emph{surjective} homomorphism
    \[
        \functiondomain{\phi}{G}{H},
    \]
    and we claim that this is injective. Indeed, take \(g \in \ker\pqty{\phi}\), then \(\phi\pqty{g} = \identity_{X}\), and therefore
    \[
        g \gamma = \gamma
    \]
    for all \(\gamma \in G\). In particular, taking \(\gamma = e \in G\) gives \(g = e\).

    Therefore, the kernel is trivial, and \(\phi\) is an injection. Therefore, \(\phi\) is a bijection,
    \[
        G \isomorphic H \subgroup \Sym\pqty{X},
    \]
    as required.

    Furthermore, notice \(X = G\), so \(\abs{G} < \infty\) allows us to choose \(\abs{X} < \infty\).
\end{proof}

\subsection{Orbits and Stabilisers}
\label{sec:orbits-and-stabilisers}

We in fact can learn a huge amount about groups by thinking about group actions.

\begin{definition}[Orbits and Stabilisers]
    \label{def:orbit-stabiliser}%
    Let \(G \actson X\) and let \(x \in X\).
    \begin{enumerate}
        \item The \emph{orbit} of \(x\) is the set
        \[
            \Orb_{G}\pqty{x} = Gx = \set{y \in X}{\exists g \in G: y = gx}.
        \]
        \item The \emph{stabiliser} of \(x\) is the set
        \[
            \Stab_{G}\pqty{x} = \set{g \in G}{g\pqty{x} = x}.
        \]
    \end{enumerate}

    Note some authors use \(G_x\) for stabilisers, where \(x\) is a subscript.

    \begin{itemize}
        \item If \(\Orb_{G}\pqty{x} = X\), we say the action is \emph{transitive}. (We will see soon that this statement is equivalent with for \emph{all} \(x \in X\) or for \emph{some} \(x \in X\).)
        \item If every element \(g \in G\) (except the identity \(g = e\)) has an \(x \in X\) such that \(gx \neq x\), then we say the action is \emph{faithful}.
    \end{itemize}
\end{definition}

\begin{remark}
    An action \(G \actson X\) is faithful is equivalent to that the associated homomorphism has a trivial kernel/is injective.
\end{remark}

\begin{proposition}[Stabilisers are Subgroups, Orbits Partition]
    \label{prop:stabilisers-subgroups-orbits-partition}%
    Suppose \(G \actson X\). Then
    \begin{enumerate}
        \item \(\Stab_{G} \pqty{x} \subgroup G\) for any \(x \in X\).
        \item The orbits
        \[
            \set{\Orb_{G}\pqty{y}}{y \in X}
        \]
        form a partition of \(X\).
    \end{enumerate}
\end{proposition}

\begin{remark}
    The second part of \autoref{prop:stabilisers-subgroups-orbits-partition} tells us that if \(\Orb\pqty{x} = X\) for some \(x \in X\), then there is only one orbit, and hence \(\Orb\pqty{x} = X\) for all \(x \in X\).
\end{remark}

\begin{proof}
    \begin{enumerate}
        \item We check that \(\Stab_{G} \pqty{x}\) satisfies \autoref{def:subgroup} (Subgroup). If \(g, h \in \Stab_{G}\pqty{x}\), then
        \begin{align*}
            \pqty{gh} x &= g\pqty{h\pqty{x}} && \text{\autoref{def:group-action} (Group Action)} \\
            &= g\pqty{x} && h \in \Stab_{G} \pqty{x} \\
            &= x, && g \in \Stab_{G} \pqty{x}
        \end{align*}
        so \(gh \in \Stab_{G} \pqty{x}\).

        Also, \(ex = x\) from \autoref{def:group-action} (the definition of a group action), so \(e \in \Stab_G \pqty{x}\).

        Finally, for \(g \in \Stab_G \pqty{x}\),
        \begin{align*}
            g\inverse x &= g\inverse \pqty{gx} && g \in \Stab_{G} \pqty{x} \\
            &= \pqty{g\inverse g} x && \text{\autoref{def:group-action} (Group Action)} \\
            &= ex \\
            &= x, && \text{\autoref{def:group-action} (Group Action)}
        \end{align*}
        so \(g\inverse \in \Stab_{G} \pqty{x}\).

        Hence, \(\Stab_{G} \pqty{x} \subgroup G\).

        \item The proof is similar to the proof of \autoref{lem:cosets-partition} (Cosets Partition).
        \begin{itemize}
            \item For any \(x \in X\), \(x = ex\) for \(e \in G\), so \(x \in \Orb_G \pqty{x}\), so indeed orbits cover \(X\).
            \item If \(\Orb_{G} \pqty{x_1} \cap \Orb_{G} \pqty{x_2} \neq \emptyset\), then there is \(y = g_1 x_1 = g_2 x_2\) for some \(g_1, g_2 \in G\).
        
            Hence,
            \begin{align*}
                x_1 &= \pqty{g_1 \inverse g_1} x_1 \\
                &= g_1 \inverse \pqty{g_1 x_1} \\
                &= g_1 \inverse \pqty{g_2 x_2} \\
                &= \pqty{g_1 \inverse g_2} x_2 \\
                &\in \Orb_{G} \pqty{x_2}.
            \end{align*}

            Moreover, any \(g x_1\) can be written as
            \[
                g x_1 = g \pqty{g_1\inverse g_2} x_2 \in \Orb_{G} \pqty{x_2}.
            \]

            Hence, \(\Orb_{G} \pqty{x_1} \subseteq \Orb_{G} \pqty{x_2}\). Likewise, \(\Orb_{G} \pqty{x_2} \subseteq \Orb_{G} \pqty{x_1}\), so we have
            \[
                \Orb_{G} \pqty{x_1} = \Orb_{G} \pqty{x_2}
            \]
            as desired.
        \end{itemize}
    \end{enumerate}
\end{proof}

\begin{example*}[\(D_{2n} \actson X_n\), the regular \(n\)-gon]
    As in \autoref{fig:dihedral-group-action-hexagon} for the hexagon case, for any \(j\),
    \begin{align*}
        r^j s \pqty{x} &= r^j \pqty{x} \\
        &= e^{\frac{2 \pi i j}{n}} \cdot 1 \\
        &= e^{\frac{2 \pi i j}{n}},
    \end{align*}
    so
    \begin{align*}
        \Orb_{D_{2n}}\pqty{x} &= \set{e^{\frac{2 \pi i j}{n}}}{0 \leq j < n} \\
        &= \Bqty{n\text{-th roots of unity}}.
    \end{align*}

    The above calculation also shows that
    \[
        gx = x \iff g \in \Bqty{r, s},
    \]
    so
    \[
        \Stab_{D_{2n}}\pqty{x} = \Bqty{r, s}.
    \]
\end{example*}

\begin{figure}[ht!]
    \centering
    \caption{\(D_{2n} \actson X_n\) for \(n = 6\), the regular hexagon}
    \label{fig:dihedral-group-action-hexagon}
\end{figure}

\begin{theorem}[Orbit--Stabiliser]
    \label{thm:orbit-stabiliser}%
    Suppose \(G \actson X\). Let \(x \in X\), the map
    \[
        g \Stab_{G}\pqty{x} \mapsto gx
    \]
    defines a well-defined bijection
    \[
        \flatfrac{G}{\Stab_{G}\pqty{x}} \to \Orb_{G}\pqty{x}.
    \]
\end{theorem}

We first prove a corollary of this theorem, naming the counting version of this.

\begin{corollary}[Orbit--Stabiliser]
    \label{cor:orbit-stabiliser}%
    If \(G \actson X\) and \(x \in X\), then
    \[
        \abs{G} = \abs{\Orb_{G}\pqty{x}} \abs{\Stab_{G}\pqty{x}}.
    \]
\end{corollary}

\begin{proof}
    \autoref{thm:orbit-stabiliser} (Orbit--Stabiliser Theorem) gives
    \[
        \abs{\Orb_{G} \pqty{x}} = \abs{\flatfrac{G}{\Stab_{G}} \pqty{x}},
    \]
    and \autoref{thm:lagrange} (Lagrange's Theorem) gives
    \[
        \abs{\flatfrac{G}{\Stab_{G} \pqty{x}}} = \frac{\abs{G}}{\abs{\Stab_{G} \pqty{x}}}.
    \]

    This gives the result as desired.
\end{proof}

Before we prove \autoref{thm:orbit-stabiliser} (Orbit--Stabiliser Theorem), we first give an example of the counting version of this.

\begin{example*}[\(D_{2n} \actson X_n\)]
    We see that if \(x \in X_n\) is a vertex, then
    \[
        \abs{\Orb_{D_{2n}} \pqty{x}} = n
    \]
    and
    \[
        \abs{\Stab_{D_{2n}} \pqty{x}} = 2.
    \]

    This gives an easy (but circular) `proof' that
    \[
        \abs{D_{2n}} = 2n.
    \]
\end{example*}

\begin{proof}[of \autoref{thm:orbit-stabiliser} (Orbit--Stabiliser Theorem)]
    We write
    \[
        S = \Stab_{G} \pqty{x},
    \]
    and define
    \[
        \phi \pqty{gS} = gx.
    \]

    We check that \(\phi\) is a well-defined bijection.

    \begin{itemize}
        \item \textbf{\(\phi\) is well-defined.} Let \(g_1, g_2 \in G\) such that \(g_1 S = g_2 S\). We need to check that \(\phi\pqty{g_1 S} = \phi\pqty{g_2 S}\). Now, \(g_1 S = g_2 S\) means that
        \[
            g_1 = g_2 s
        \]
        for some \(s \in S\).

        Hence,
        \begin{align*}
            g_1 x &= \pqty{g_2 s} x \\
            &= g_2 \pqty{sx} \\
            &= g_2 x
        \end{align*}
        since \(s \in S\) gives \(sx = x\), as required.

        \item \textbf{\(\phi\) is surjective.} For any \(gx \in \Orb_{G} \pqty{x}\), \(g \in G\), we note
        \[
            \phi\pqty{gS} = gx
        \]
        so \(\phi\) is surjective.

        \item \textbf{\(\phi\) is injective.} Suppose that \(\phi\pqty{g_1 S} = \phi\pqty{g_2 S}\), which is equivalent to \(g_1 x = g_2 x\). We want to show \(g_1 S = g_2 S\). Indeed, let \(g_1 = g_2 s\) for some \(s\), we want to show \(s \in S\):
        \begin{align*}
            sx &= \pqty{g_2 \inverse g_1} x \\
            &= g_2 \inverse \pqty{g_1 x} \\
            &= g_2 \inverse \pqty{g_2 x} \\
            &= \pqty{g_2 \inverse g_2} x \\
            &= ex \\
            &= x,
        \end{align*}
        so \(s \in \Stab_{G} \pqty{x}\), and hence \(s \in S\).

        Therefore,
        \begin{align*}
            g_1 &= \pqty{g_2 g_2\inverse} g_1 \\
            &= g_2 \pqty{g_2\inverse g_1} \\
            &= g_2 s \\
            &= g_2 S,
        \end{align*}
        and since cosets partition, it follows that \(g_1 S = g_2 S\) as required.
    \end{itemize}
\end{proof}

\begin{example*}[Symmetries of a Cube]
    Let \(G\) be the group of isometries of a cube.

    As in \autoref{fig:symmetries-of-a-cube}, let \(x\) be the centre of a face of a cube. Then
    \[
        \abs{\Orb_{G} \pqty{x}} = 6,
    \]
    while
    \[
        \Stab_{G}\pqty{x} \isomorphic D_8,
    \]
    so
    \[
       \abs{\Stab_{G}\pqty{x}} = 8.
    \]

    Therefore,
    \begin{align*}
        \abs{G} &= \abs{\Orb_{G} \pqty{x}} \abs{\Stab_{G} \pqty{x}} \\
        &= 6 \cdot 8 \\
        &= 48.
    \end{align*}
\end{example*}

\begin{figure}[ht!]
    \centering
    \caption{A Cube and the Centre of a Face}
    \label{fig:symmetries-of-a-cube}
\end{figure}

The next theorem is a different kind of application of \autoref{cor:orbit-stabiliser} (Orbit--Stabiliser Theorem). The proof is not very intuitive, but gives a very important result in group theory.

\begin{theorem}[Cauchy's Theorem]
    \label{thm:cauchy}
    If \(\abs{G} < \infty\), and \(p\) is a prime that divides \(\abs{G}\), then there is a \(g \in G\) such that
    \[
        \ord\pqty{g} = p.
    \]
\end{theorem}

\begin{proof}
    Consider the set \(X\) of \(p\)-tuples
    \[
        X = \set{\pqty{g_1, \ldots, g_p}}{g_i \in G, g_1 g_2 \cdots g_p = e}.
    \]

    We define an action of \(C_p\) on \(X\), where
    \[
        C_p = \Bqty{1, t, t^2, \ldots, t^{p - 1}} = \angBrac{t}
    \]
    as follows. Let
    \[
        t^k \pqty{g_1, \ldots, g_p} = \pqty{g_{k + 1}, \ldots, g_p, g_1, \ldots, g_k}.
    \]

    It is easy to see that
    \[
        \pqty{t^k \cdot t^l} \pqty{g_1, \ldots, g_p} = \pqty{t^k} \pqty{t^l \pqty{g_1, \ldots, g_p}}.
    \]

    We also need to check that
    \[
        \pqty{g_{k + 1}, \ldots, g_p, g_1, \ldots, g_k} \in X.
    \]

    Suppose that
    \[
        a = g_1 \cdots g_k, b = g_{k + 1} \cdots g_p,
    \]
    we have \(ab = e\), and hence \(b = a\inverse\), therefore
    \[
        ba = g_{k + 1} \cdots g_p g_1 \cdots g_k = e
    \]
    as required. This means we indeed defined an action \(C_p \actson X\).

    We now compute the size of \(X\). We may decide \(g_1, \ldots, g_{p - 1} \in G\) to be arbitrary elements, then \(g_p\) has to be determined uniquely, namely
    \[
        g_p = g_{p - 1}\inverse \cdots g_2\inverse g_1\inverse
    \]
    such that
    \[
        \pqty{g_1, \ldots, g_{p - 1}, g_p} \in X.
    \]

    Therefore,
    \[
        \abs{X} = \abs{G}^{p - 1}.
    \]

    From \autoref{prop:stabilisers-subgroups-orbits-partition}, \(C_p \actson X\) partition \(X\) into orbits
    \[
        X = \Orb\pqty{x_1} \cup \Orb\pqty{x_2} \cup \cdots \cup \Orb\pqty{x_k}.
    \]

    Hence, by \autoref{cor:orbit-stabiliser} (Orbit--Stabiliser Theorem), for each \(j = 1, 2, \ldots, k\),
    \[
        \abs{\Orb\pqty{x_j}} \divides \abs{C_p} = p
    \]
    so either the size of \(\Orb\pqty{x_j}\) is \(1\), or \(p\).

    Let \(l\) be the number of orbits with size \(1\). After re-numbering, we may assume that
    \[
        \abs{\Orb\pqty{x_j}} = \begin{cases}
            1, & 1 \leq j \leq l, \\
            p, & j > l.
        \end{cases}
    \]

    Since the orbits partition \(X\),
    \begin{align*}
        \abs{G}^{p - 1} &= \abs{X} \\
        &= \sum_{j = 1}^{k} \abs{\Orb\pqty{x_j}} \\
        &= \sum_{j = 1}^{l} \abs{\Orb{\pqty{x_j}}} + \sum_{j = l + 1}^{k} \abs{\Orb{\pqty{x_j}}} \\
        &= l + p \cdot \pqty{k - l}.
    \end{align*}

    Since \(p \divides \abs{G}\), it follows that \(p \divides l\).

    Now, note that \(x = \pqty{g_1, \ldots, g_p}\) has an orbit of size \(1\) if and only if \(g_1 = g_2 = \cdots = g_p\).

    In particular, there is at least one such element giving such an orbit, namely
    \[
        \pqty{e, e, \ldots, e} \in X.
    \]

    So \(l > 0\). Since \(p \divides l\), it follows that
    \[
        l \geq p > 1,
    \]
    so there is at least \(1\) more orbit
    \[
        \pqty{g, g, \ldots, g} \in X
    \]
    with \(g \neq e\).

    The definition of \(X\) implies that \(g^p = e\), and hence we can see \(\ord\pqty{g} \divides p\), and hence \(\ord\pqty{g} = p\) since \(p\) is prime and \(g \neq e\), as required.
\end{proof}

This is a proof that is not very intuitive, but the result is very useful.

\subsection{Conjugation}
\label{sec:conjugation}

\begin{definition}[Conjugate, Conjugacy Class]
    \label{def:conjugate-ccl}%
    Let \(G\) be a group, and \(g \in G\). For any \(h \in G\), the element \(hgh\inverse \in G\) is called \emph{the conjugate of \(g\) by \(h\)}.

    The \emph{conjugacy class} of \(g \in G\) is
    \[
        \ccl_{G}\pqty{g} = \set{hgh\inverse}{h \in G}.
    \]
\end{definition}

The intuition behind this definition is that \(hgh\inverse\) has the same `shape' as \(g\), and \(hgh\inverse\) corresponds to a 'change in coordinates' of \(g\) by \(h\). Conjugacy classes partition the group.

\begin{example*}[Conjugacy in an Abelian Group]
    If \(G\) is abelian, then for any \(g, h \in G\),
    \[
        hgh\inverse = ghh\inverse = ge = g,
    \]
    so the only element conjugate to \(g\) is \(g\) itself. All conjugacy classes have size \(1\). (This condition is equivalent to a group being abelian.)
\end{example*}

Conjugacy classes of reflections in a dihedral group demonstrates the intuition of `shape' well.

\begin{example*}[Conjugacy Classes of Reflections in Dihedral Group]
    In Sheet 2, Q6, \(D_{2n}\) has \(1\) conjugacy class of reflections if \(n\) is odd, and \(2\) if \(n\) is even.
    
    As in \autoref{fig:ccl-dihedral}, when \(n = 6\), there are two `types' of reflections, those whose axis run through vertices, and those who run through midpoints of edges. When \(n = 5\), there is only one `type' of reflection: their axis all run through a vertex and a midpoint of an edge.
\end{example*}

\begin{figure}[ht!]
    \centering
    \caption{Conjugacy in Dihedral Groups}
    \label{fig:ccl-dihedral}
\end{figure}

Conjugation defines another action of \(G \actson G\):
\[
    g \star \gamma = g \gamma g\inverse
\]
for \(g \in G, \gamma \in X = G\). This is \textbf{very different} to the left regular action. In terms of language in actions,
\[
    \ccl_{G} \pqty{g} = \Orb_{G\star} \pqty{g}.
\]

From the example on abelian groups, we see that conjugating to an element itself is closely related to commutativity:
\[
    hgh\inverse = g \iff hg = gh.
\]

This leads to \autoref{def:centraliser-centre}.

\begin{definition}[Centraliser, Centre]
    \label{def:centraliser-centre}%
    The \emph{centraliser} of \(g \in G\) is
    \[
        C_G\pqty{g} = \set{h \in G}{h g h\inverse = g}.
    \]

    It is the elements that \emph{commute with \(g\)}.

    The \emph{centre} of \(G\) is
    \[
        Z\pqty{G} = \set{h \in G}{\forall g \in G, hg = gh} = \bigcap_{g \in G} C_G \pqty{g}.
    \]

    It is the elements that \emph{commute with everything in \(G\)}.
\end{definition}

Using the language of the conjugation action, we have
\[
    C_G\pqty{g} = \Stab_{G\star} \pqty{g}
\]
and hence \(C_G\pqty{g} \subgroup G\) forms a group. Since \(Z\pqty{G}\) is an intersection of subgroups, it is itself a subgroup: \(Z\pqty{G} \subgroup G\).

\begin{remark}
    Another useful way to think of the centre is
    \[
        Z\pqty{G} = \bigcup_{g \in G, \abs{\ccl\pqty{g}} = 1} \ccl\pqty{g} = \bigcup_{g \in G, \abs{\ccl\pqty{g}} = 1} \Bqty{g},
    \]
    i.e. the union of the conjugacy classes with size one. (We will see later that this means \(Z\pqty{G} \normalsub G\).)
\end{remark}